% Esqueleto basado en los informes IEEE de micro/Laboratorios.
% Completar autores, textos entre corchetes y evidencias antes de entregar.
% La consigna limita el informe completo a 6 paginas.
% Archivo autocontenido para el editor de Codex u Overleaf.
\documentclass[conference]{IEEEtran}
\usepackage[T1]{fontenc}
\usepackage[utf8]{inputenc}
\usepackage[spanish,es-tabla]{babel}
\usepackage{amsmath,amssymb}
\usepackage{graphicx}
\usepackage{booktabs}
\usepackage{xcolor}
\usepackage{listings}
\usepackage[hidelinks]{hyperref}
\usepackage[left=2cm,right=2cm,top=2cm,bottom=2cm]{geometry}

\lstset{language=VHDL,basicstyle=\ttfamily\footnotesize,
  frame=single,breaklines=true,columns=fullflexible,
  keepspaces=true,showstringspaces=false,captionpos=b}
\newcommand{\pendiente}[1]{\textit{[Completar: #1.]}}
\newcommand{\evidencia}[1]{%
  \fbox{\parbox[c][2.4cm][c]{0.90\columnwidth}{\centering\small #1}}}

\begin{document}

\title{Electrónica Digital 2: Laboratorio 1\\
Introducción a Vivado y a la placa Spartan Edge Accelerator}
\author{\IEEEauthorblockN{[Nombre y apellido]}
\IEEEauthorblockA{\textit{Ingeniería en Mecatrónica}\\
\textit{Universidad Tecnológica (UTEC)}\\
Fray Bentos, Uruguay\\
{[Correo institucional]}}
\and
\IEEEauthorblockN{[Nombre y apellido]}
\IEEEauthorblockA{\textit{Ingeniería en Mecatrónica}\\
\textit{Universidad Tecnológica (UTEC)}\\
Fray Bentos, Uruguay\\
{[Correo institucional]}}}
% Agregar un tercer bloque con \and si corresponde; maximo 3 integrantes.
\maketitle

\begin{abstract}
\pendiente{redactar al finalizar el trabajo un resumen del objetivo,
del circuito combinacional desarrollado en VHDL, de su verificación
mediante testbench y de su implementación en la FPGA. Indicar el
resultado principal observado sin introducir datos no medidos}
\end{abstract}
\begin{IEEEkeywords}
FPGA, VHDL, Vivado, Spartan-7, lógica combinacional, testbench.
\end{IEEEkeywords}

\section{Introducción}
El diseño digital sobre una FPGA permite convertir una descripción
de hardware en un circuito configurable. Este laboratorio de Electrónica
Digital 2 propone una primera aproximación a ese proceso utilizando
VHDL, Vivado y la placa Spartan Edge Accelerator con el shield UTEC.
El objetivo es relacionar la especificación lógica de un circuito con
su simulación y, posteriormente, con su comportamiento físico.

La actividad comienza con una compuerta AND de dos entradas para
familiarizarse con la herramienta. El ejercicio evaluado consiste en
un circuito combinacional de tres entradas ($A$, $B$ y $C$) y tres
salidas ($X$, $Y$ y $Z$), definido por una tabla de verdad. Su
verificación requiere recorrer las ocho combinaciones tanto en el
testbench como en la placa, mediante dip-switches y LEDs~\cite{guia}.

\subsection{Objetivos}
\begin{itemize}
\item Comprender las etapas del flujo de diseño en FPGA con Vivado.
\item Distinguir la entidad y la arquitectura de una descripción VHDL.
\item Describir un circuito combinacional a partir de su tabla de verdad.
\item Verificar las ocho combinaciones mediante simulación funcional.
\item Asignar las señales a los pines y validar el circuito en la placa.
\end{itemize}

\section{Marco teórico}
\subsection{FPGA y flujo de diseño}
Una FPGA contiene recursos lógicos e interconexiones configurables.
Para un circuito combinacional, las salidas dependen de las entradas
actuales, sin requerir almacenamiento del estado anterior. La
descripción VHDL expresa relaciones de hardware que pueden operar
concurrentemente; no representa una secuencia de instrucciones
ejecutadas por un microcontrolador.

El flujo parte de la descripción y su simulación funcional. La síntesis
transforma el código RTL en una red lógica; la implementación ubica
y conecta los recursos en la FPGA. Luego se genera el bitstream que
configura el dispositivo~\cite{vivado_impl}. En este laboratorio se
prevé cargarlo por JTAG desde Vivado~\cite{guia}.

\subsection{Descripción de circuitos en VHDL}
La entidad define la interfaz externa, incluidos los nombres, modos
y tipos de los puertos. La arquitectura describe el comportamiento o
la estructura que corresponde a esa interfaz~\cite{vhdl_entidad}.
Se propone agrupar las entradas como $ABC$, con $A$ como bit más
significativo, y las salidas como $XYZ$, con $X$ en esa misma posición.

Una asignación seleccionada \texttt{with/select} permite asociar cada
entrada con una salida; \texttt{when/else} expresa alternativas
condicionales. La descripción debe asignar las salidas para todos
los casos, evitando almacenamiento involuntario. El tipo
\texttt{STD\_LOGIC} incluye valores distintos de \texttt{'0'} y
\texttt{'1'}, como \texttt{'U'} y \texttt{'X'}; por ello,
\texttt{others} permite definir el comportamiento del modelo ante
combinaciones no binarias. Ese tratamiento no reemplaza la
verificación de las ocho filas de la tabla.

\subsection{Testbench y simulación}
El testbench instancia el circuito y aplica estímulos sin formar parte
del hardware que se programa. Para tres entradas binarias existen
$2^3=8$ combinaciones. Se prevé aplicar cada una, esperar la
actualización de las señales y comparar las salidas con la tabla.
La simulación RTL comprueba el modelo; una simulación temporal de
un diseño implementado puede incorporar sus retardos~\cite{simulacion}.

\subsection{Placa SEA y shield UTEC}
La Spartan Edge Accelerator integra una FPGA Spartan-7 XC7S15 y
admite conexión con shields de formato Arduino~\cite{sea}.
Según la guía, el shield UTEC proporciona los dip-switches y LEDs
que se usarán como entradas y salidas. Las señales se configurarán
con \texttt{LVCMOS33}. La alimentación de 3,3 V de los dip-switches
requiere llevar \texttt{AR\_3V3\_EN}, conectado al pin L13, a
\texttt{'0'}~\cite{guia}. La correspondencia de las otras seis señales
con pines físicos queda pendiente de consultar el manual del shield.

\section{Metodología}
\subsection{Materiales y herramientas}
\begin{itemize}
\item Spartan Edge Accelerator Board y shield multifunción UTEC.
\item Cable USB-C y programador JTAG HS2 o HS3: \pendiente{modelo usado}.
\item Vivado: \pendiente{versión utilizada; la guía indica 2025.1}.
\item Manuales de la SEA y del shield disponibles en Moodle.
\end{itemize}
\pendiente{registrar el dispositivo seleccionado y confirmar con el
manual la parte \texttt{xc7s15ftgb196-1} indicada en la guía}

\subsection{Familiarización con Vivado: compuerta AND}
\pendiente{resumir la creación del proyecto RTL, la descripción de la
AND, su testbench, las restricciones y la prueba en placa. Mantener
esta parte breve para priorizar el ejercicio evaluado}

\subsection{Circuito combinacional solicitado}
La Tabla~\ref{tab:verdad} reproduce la especificación de la consigna.
Se considera a $A$ como el bit más significativo de la entrada.
\begin{table}[htbp]
\centering
\caption{Tabla de verdad del circuito solicitado}
\label{tab:verdad}
\begin{tabular}{ccc|ccc}
\toprule
$A$ & $B$ & $C$ & $X$ & $Y$ & $Z$\\
\midrule
0 & 0 & 0 & 1 & 0 & 0\\
0 & 0 & 1 & 0 & 1 & 1\\
0 & 1 & 0 & 0 & 0 & 1\\
0 & 1 & 1 & 1 & 1 & 0\\
1 & 0 & 0 & 0 & 1 & 0\\
1 & 0 & 1 & 1 & 1 & 1\\
1 & 1 & 0 & 0 & 0 & 0\\
1 & 1 & 1 & 1 & 0 & 1\\
\bottomrule
\end{tabular}
\end{table}
\pendiente{explicar los puertos, los vectores de entrada y salida, la
sentencia elegida y el caso \texttt{others}. Incorporar un fragmento
representativo del código y remitir al archivo VHDL completo}

\subsection{Testbench y simulación funcional}
\pendiente{describir la instancia del circuito, el tiempo entre
estímulos, el recorrido de las ocho combinaciones y el criterio de
comparación con la tabla. Indicar la duración total de la simulación}

\subsection{Restricciones y programación de la placa}
\pendiente{completar la Tabla~\ref{tab:pines} a partir del manual del
shield. Explicar síntesis, implementación, revisión de mensajes,
generación del bitstream y programación por Hardware Manager}
\begin{table}[htbp]
\centering
\caption{Asignación de señales en el archivo XDC}
\label{tab:pines}
\small
\begin{tabular}{lll}
\toprule
Señal & Elemento físico & Pin FPGA\\
\midrule
$A$ & [Dip-switch] & [Pin]\\
$B$ & [Dip-switch] & [Pin]\\
$C$ & [Dip-switch] & [Pin]\\
$X$ & [LED] & [Pin]\\
$Y$ & [LED] & [Pin]\\
$Z$ & [LED] & [Pin]\\
\texttt{AR\_3V3\_EN} & Habilitación 3,3 V & L13\\
\bottomrule
\end{tabular}
\end{table}
Según la guía, se utiliza \texttt{LVCMOS33} y se fuerza
\texttt{AR\_3V3\_EN} a \texttt{'0'} para habilitar la alimentación
de los dip-switches.

\section{Resultados y discusión}
\subsection{Resultados de simulación}
\pendiente{incluir y analizar la captura de las ocho combinaciones;
identificar entradas, salidas, escala temporal y posibles discrepancias}
\begin{figure}[htbp]
\centering
\evidencia{Insertar captura de la simulación en Vivado.\\
Deben verse $A$, $B$, $C$, $X$, $Y$ y $Z$.}
% Reemplazar \evidencia por \includegraphics al incorporar la captura.
\caption{Espacio reservado para la verificación por simulación.}
\label{fig:simulacion}
\end{figure}

\subsection{Validación en hardware}
\pendiente{documentar las pruebas en placa y completar las columnas
de observaciones de la Tabla~\ref{tab:validacion}. No marcar coincidencia
hasta realizar las pruebas}
\begin{table}[htbp]
\centering
\caption{Comparación entre especificación, simulación y placa}
\label{tab:validacion}
\begin{tabular}{cccc}
\toprule
$ABC$ & $XYZ$ esperado & Simulado & En placa\\
\midrule
000 & 100 & [---] & [---]\\
001 & 011 & [---] & [---]\\
010 & 001 & [---] & [---]\\
011 & 110 & [---] & [---]\\
100 & 010 & [---] & [---]\\
101 & 111 & [---] & [---]\\
110 & 000 & [---] & [---]\\
111 & 101 & [---] & [---]\\
\bottomrule
\end{tabular}
\end{table}
\begin{figure}[htbp]
\centering
\evidencia{Insertar fotografía de la SEA con el shield.\\
Identificar los tres dip-switches y los tres LEDs.}
\caption{Espacio reservado para el montaje utilizado.}
\label{fig:montaje}
\end{figure}
\pendiente{agregar enlace al video que muestre las ocho combinaciones,
e indicar incidencias y cómo se resolvieron}

\section{Preguntas del laboratorio}
\subsection{Entidad y arquitectura}
La entidad establece qué señales entran y salen del circuito; la
arquitectura especifica cómo se relacionan. Una entidad puede tener
varias arquitecturas, permitiendo describir implementaciones
alternativas de una misma interfaz~\cite{vhdl_entidad}. Para este
ejercicio, una arquitectura podría describir la tabla y otra las
ecuaciones booleanas, conservando los mismos puertos.

\subsection{Función del archivo XDC}
El archivo XDC reúne restricciones del diseño. En este laboratorio,
\texttt{PACKAGE\_PIN} vincula los puertos con terminales físicos y
\texttt{IOSTANDARD} define el estándar eléctrico. También existen
restricciones temporales~\cite{xdc}. Sin la asignación de pines no
queda establecida la conexión con los dip-switches y LEDs. La lógica
puede simularse, pero no se garantiza el conexionado previsto en
hardware; las comprobaciones de Vivado pueden impedir generar el
bitstream mientras existan puertos sin restricciones adecuadas.

\subsection{Estándar LVCMOS33}
\texttt{LVCMOS33} corresponde a señales CMOS referidas a una
alimentación nominal de 3,3 V. La guía exige ese estándar para las
conexiones utilizadas~\cite{guia}. Elegir un estándar en el XDC no
convierte físicamente una señal de 5 V en una de 3,3 V: los niveles
aplicados deben ser compatibles con el banco de E/S y sus límites
eléctricos. Por ello no corresponde asumir compatibilidad directa
con 5 V. La identificación del banco y sus límites concretos se
confirmará con la documentación de la placa.

\subsection{Habilitación del riel de 3,3 V}
Según la consigna, \texttt{AR\_3V3\_EN} controla un interruptor de
alimentación mediante MOSFET, activo en bajo. Debe declararse como
salida, asignarse al pin L13 y mantenerse en \texttt{'0'} para
habilitar el riel del conector Arduino. Sin esa habilitación, los
dip-switches carecen de alimentación y, según la guía, sus entradas
se leen como cero. Los LEDs conectados al header FPGAIO no dependen
de esa habilitación~\cite{guia}. Esta explicación aún no constituye
una observación experimental del montaje.

\subsection{Expresión booleana mínima}
Se elige la salida $X$, cuyos unos aparecen en las entradas $000$,
$011$, $101$ y $111$. Con $A$ como bit más significativo:
\begin{equation}
X=\sum m(0,3,5,7).
\end{equation}
En un mapa de Karnaugh, los minterms $3$ y $7$ se agrupan para
obtener $BC$, y los minterms $5$ y $7$ producen $AC$. El minterm
$0$ queda aislado, pues todas sus celdas adyacentes valen cero.
La suma de productos mínima resulta:
\begin{equation}
\boxed{X=\overline{A}\,\overline{B}\,\overline{C}+BC+AC.}
\label{eq:x_minima}
\end{equation}
El término del minterm $0$ es indispensable; $BC$ es necesario para
cubrir $3$ y $AC$ para cubrir $5$. Al evaluar la ecuación en orden
de $000$ a $111$, se obtiene $1,0,0,1,0,1,0,1$, coincidiendo con
la columna $X$ de la Tabla~\ref{tab:verdad}. Esta equivalencia se
refiere a entradas binarias estables; no demuestra igualdad de
retardos físicos ni determina el tratamiento de valores no binarios.

\subsection{Simulación funcional y placa real}
La simulación funcional evalúa la descripción RTL con los estímulos
del testbench. La placa agrega retardos de propagación, condiciones
eléctricas y elementos físicos. Por ejemplo, el rebote de un
dip-switch o una alimentación deshabilitada no aparecen en un
testbench que solo aplica cambios binarios ideales. Al cambiar
varios interruptores también pueden producirse combinaciones
intermedias. Por tanto, una coincidencia en simulación debe
complementarse con la prueba física. Vivado distingue simulación
RTL de simulación posterior a implementación con información
temporal~\cite{simulacion}.

\section{Conclusiones}
\pendiente{evaluar el cumplimiento de los objetivos con base en los
resultados, sintetizar lo aprendido sobre el flujo FPGA y describir
limitaciones o mejoras. Evitar afirmar éxito antes de verificarlo}

\begin{thebibliography}{00}
\bibitem{guia} UTEC, ``Laboratorio 1: Introducción a Vivado y a la placa
Spartan Edge Accelerator,'' Electrónica Digital 2, guía de laboratorio.
% Agregar solamente fuentes efectivamente consultadas y citarlas en el texto.
% Completar datos de los manuales SEA/shield y documentación VHDL/Vivado.
\bibitem{vivado_impl} AMD, ``Design Entry and Implementation,'' [En línea]. Disponible: \url{https://www.amd.com/en/products/software/adaptive-socs-and-fpgas/vivado/implementation.html}. Consultado: 2 de octubre de 2026.
\bibitem{vhdl_entidad} AMD, ``VHDL Entity and Architecture Descriptions, UG901,'' [En línea]. Disponible: \url{https://docs.amd.com/r/en-US/ug901-vivado-synthesis/VHDL-Entity-and-Architecture-Descriptions}. Consultado: 2 de octubre de 2026.
\bibitem{xdc} AMD, ``Vivado Design Suite User Guide: Using Constraints, UG903,'' [En línea]. Disponible: \url{https://docs.amd.com/r/en-US/ug903-vivado-using-constraints}. Consultado: 2 de octubre de 2026.
\bibitem{simulacion} AMD, ``Vivado Design Suite User Guide: Logic Simulation, UG900,'' [En línea]. Disponible: \url{https://docs.amd.com/r/en-US/ug900-vivado-logic-simulation}. Consultado: 2 de octubre de 2026.
\bibitem{sea} Seeed Studio, ``Spartan Edge Accelerator Board,'' [En línea]. Disponible: \url{https://wiki.seeedstudio.com/Spartan-Edge-Accelerator-Board/}. Consultado: 2 de octubre de 2026.
\end{thebibliography}

% Entregables externos al informe: VHDL completo, testbench, captura,
% archivo XDC y video de las ocho combinaciones.
% Evitar anexar codigo extenso si compromete el maximo de 6 paginas.
\end{document}
